%%=============================================================================
%% DADOS DO TRABALHO -- FONTE ÚNICA DA VERDADE
%%
%% Este é o ÚNICO arquivo onde os dados do seu trabalho devem ser digitados.
%% Tudo o que aparece na capa, folha de rosto, ficha catalográfica, folha de
%% aprovação e nos metadados do PDF é gerado a partir daqui.
%%
%% Campos marcados como OBRIGATÓRIO que ficarem em branco aparecem no PDF
%% como "[[ PREENCHER: campo ]]" e geram aviso na compilação. Ao gerar a
%% versão final, use a opção `entrega' em main.tex: ela transforma o aviso
%% em erro e impede a compilação com campos pendentes.
%%
%% Para reaproveitar estes dados no corpo do texto, use os acessores:
%%   \TituloTrabalho  \NomeAutor      \NomeOrientador   \NomeCoorientador
%%   \NomeCurso       \NomePrograma   \NomeInstituicao  \AreaConcentracao
%%   \LinhaPesquisa  \LocalTrabalho   \CidadeTrabalho   \AnoTrabalho
%% Exemplo: "Ao meu orientador, \NomeOrientador, agradeço..."
%%=============================================================================

%% --- Identificação ---------------------------------------------------------
%% Título de trabalho, proposto. Ajuste livremente.
\titulo{Exposição ocupacional à inteligência artificial e transições no
        mercado de trabalho brasileiro}                    % OBRIGATÓRIO
\subtitulo{evidências de painel da PNAD Contínua após o lançamento
           público do ChatGPT}                             % opcional
\titulotraduzido{Occupational exposure to artificial intelligence and labor
                 market transitions in Brazil: panel evidence from the
                 Brazilian Continuous National Household Sample Survey after
                 the public release of ChatGPT}             % opcional

\autor{Ricardo Carvalho}                                   % OBRIGATÓRIO
\autorficha{CARVALHO, Ricardo}                             % OBRIGATÓRIO

\orientador{Prof. Dr. Danny de Castro Soares}              % OBRIGATÓRIO
\coorientador{}                                            % opcional

%% --- Vínculo institucional -------------------------------------------------
\curso{}                                 % OBRIGATÓRIO em bacharelado/especialização
\programa{Administração Pública}         % OBRIGATÓRIO em mestrado/doutorado
\habilitacao{Mestrado Profissional em Administração Pública}  % linha da capa
\area{Ciência de Dados e Inteligência Artificial}  % área de concentração (pós-graduação)
\linhapesquisa{}                         % linha de pesquisa (pós-graduação)

\cidade{Brasília}                        % usada na folha de aprovação
\local{Brasília-DF}                      % usado na capa e folha de rosto
\data{2026}                              % ano de depósito

%% --- Defesa ----------------------------------------------------------------
%% Informe APENAS a data; a cidade vem de \cidade acima.
\datadefesa{17 de novembro de 2026}      % OBRIGATÓRIO -> "Cidade, dd de mes de aaaa."

%% Banca (modelo 2026), um comando por membro:
%%   \membrobanca[Instituição]{Nome}{Função}[ID Lattes ou ORCID]
%% Na folha de aprovação sai: traço | Nome | ID Lattes ou ORCID |
%% "Função - Instituição". A instituição vale IDP quando omitida; para o
%% examinador externo, informe a sigla da IES dele. Dado que ainda falta vai
%% em \pendente{...}: aparece como [[ PREENCHER: ... ]] e gera aviso.
%% O modelo pede um ID por membro; os três ORCID foram conferidos em orcid.org.
%% Lattes, se preferir trocar: Claudiomar 4947956711702057
%% (http://lattes.cnpq.br/4947956711702057); Rodrigo 1491651107420281
%% (http://lattes.cnpq.br/1491651107420281).
\membrobanca{Prof. Dr. Danny de Castro Soares}{Orientador}[ORCID: 0000-0002-6665-9647]
\membrobanca{Prof. Dr. Claudiomar Matias Rolim Filho}{Examinador interno}[ORCID: 0000-0003-2884-8639]
\membrobanca[UCB]{Dr. Rodrigo Batista de Araújo}{Examinador externo}[ORCID: 0009-0006-7216-2983]

%% --- Ficha catalográfica ---------------------------------------------------
%% Modelo 2026: a folha fica reservada ("Espaço destinado à ficha
%% catalográfica.") até a Biblioteca Ministro Moreira Alves enviar a ficha
%% oficial. Solicite por biblioteca@idp.edu.br ou pelo formulário da
%% Biblioteca. Quando ela chegar, salve o PDF em figuras/ e descomente:
% \arquivoficha{figuras/ficha-catalografica.pdf}

%% --- Resumo ----------------------------------------------------------------
%% NBR 6028: palavras-chave separadas por ponto e vírgula, em minúscula
%% (exceto nomes próprios), terminadas em ponto.
\palavraschave{inteligência artificial; exposição ocupacional à inteligência
               artificial; mercado de trabalho; mobilidade ocupacional;
               diferenças em diferenças; políticas públicas.}  % OBRIGATÓRIO
\keywords{artificial intelligence; occupational exposure to artificial
          intelligence; labor market; occupational mobility;
          difference-in-differences; public policy.}           % OBRIGATÓRIO
